\chapter{Defining economic variables and measuring them}
\label{exec17:ontology}

The application examples leave two connected questions: what the records
measure, and what the available evidence permits us to assert about them.
This chapter develops both. It begins with constructs, measurement error,
units and clocks; compares ontology and annotation tools; then examines
conflicting expert evidence and joint graph advice. The final Bayesian study
shows why response frequencies and accepted graph probabilities must remain
distinct. Each step asks what new information an assisted workflow supplies.

A servicing proposal may ask whether an early support offer reduces financial
distress over the next thirty days. The bank can observe payment events,
balances, customer messages and employees' notes, but none of these is simply
the construct of financial distress. A payment can be made by drawing on
another account; a note can describe anticipated hardship rather than a
realized missed payment; an employee can write more carefully after making an
offer. An assistant that assigns all these records the same variable name
changes the observation on which the causal question depends.

The proposed ontology layer has a specific purpose: make the relationship
between a construct, its operational measure and its original observation
explicit enough to challenge. This can help a scientist formulate a question,
a data engineer construct a reproducible observation and a reviewer locate a
disputed assumption. It does not determine which construct the bank should
value, establish the truth of a causal relation, or make an unobserved
confounder observable. The distinction has consequences even when the
intervention itself is randomized.

\section{From a business construct to a defined estimand}

Let \(A_i\in\{0,1\}\) denote a support offer at decision time zero for eligible
account \(i\). Define \(Y_i(a)\) as the indicator of a specified missed-payment
event by day thirty under intervention \(a\). With a fixed eligible population
and account weighting, the target is
\[
 \tau_Y=E\{Y_i(1)-Y_i(0)\}.
\]
Consistency requires a sufficiently specified offer, comparator and event;
random assignment and appropriate handling of interference would justify a
simple experimental comparison. An observational proposal additionally needs
a defensible adjustment or other identification argument. Calling a field a
confounder does not supply that argument.

Now let \(T_i(a)\) be a service note and \(L_i(a)=g_\theta(T_i(a))\) its binary
label under a frozen extraction rule \(g_\theta\). Even if the label is named
after missed payments, the experiment on labels targets
\[
 \tau_L=E\{g_\theta(T_i(1))-g_\theta(T_i(0))\}.
\]
This is a defined effect of the offer on a coded textual outcome. Equating it
with \(\tau_Y\) requires an additional measurement argument. If the desired
target is latent hardship \(U_i(a)\), neither an observed payment nor a
plausible description of it automatically identifies that target.

Egami and colleagues develop this distinction through a learned text mapping:
the representation helps define the causal outcome, and learning it from the
same treatment-assigned evaluation documents can induce dependence and
overfitting. Their train/test separation holds the learned mapping fixed for
effect estimation on the evaluation sample.~\citep{exec13_census_texts}
That remedy addresses reuse of evaluation information; it does not establish
that a learned category measures hardship without error. The earlier
measurement discussion's residual-confounding example makes a different
point: even a useful proxy can leave confounding and reverse a coefficient's
sign. Both problems survive a syntactically valid ontology.

Keith, Jensen and O'Connor explain why the causal role of text also matters:
text may measure confounding, mediate a treatment or be selected by a collider
mechanism.~\citep{exec13_census_text_review_arxiv} A note written after the offer
cannot be admitted as a baseline covariate merely because it mentions a
pre-existing condition. Its information might have been elicited by the offer.
The proposal must distinguish the underlying condition's event time, the
record's availability and the process that produced the record.

For each variable, the proposed representation retains five linked objects:
the business construct, the operational definition, the observation or
transformation, the role in a particular causal model, and the target
population's unit and weighting. These links can have different epistemic
statuses. A schema can verify that a payment record has a date and amount;
a reviewer can accept a definition for a study; a causal analyst can still
leave the identification assumptions unresolved. A single confidence score
cannot substitute for these separate decisions.

\section{What formal ontologies and language models can contribute}

Pressat-Laffouilhere's OntoBioStat thesis treats variable selection as a problem
of formal representation and supplied relations. Its account of recoding is
particularly relevant: changing a variable's categories creates a different
variable, whose context, possible values, missingness and transformation must
be represented. It also distinguishes hypothetical causal relations from
relations under validation and illustrates sensitivity analysis by retaining
or removing uncertain links before rerunning a reasoner.~\citep{exec17_ontology_thesis}
For a bank, an analogous distinction would separate a note mentioning
hardship, a customer reporting hardship and reviewed evidence supporting an
effect of hardship on payment. An inference explanation can show which
supplied relations were used. It cannot establish that those relations are
true, or guarantee a valid adjustment set for every causal graph.

The two knowledge-graph surveys offer complementary architectural choices.
Jeff Z. Pan and colleagues contrast explicit-knowledge-first with
parametric-knowledge-first approaches. The former can expose a curated
vocabulary and links for the assistant to retrieve; the latter can propose
entities and relations that subsequently need canonicalization and
review.~\citep{exec17_ontology_kg_challenges} Shirui Pan and colleagues organize
the broader interaction into knowledge-graph-enhanced language models,
language-model-augmented graphs and synergized systems. Their retrieval,
representation and graph/text fusion mechanisms explain how an assistant might
find a definition or propose a relation.~\citep{exec17_ontology_kg_roadmap}
The question-answering paths and attention scores in these architectures do
not identify the effect of an intervention.

A useful design keeps relation types separate. A type-of relation
belongs to a taxonomy; derivation describes data lineage; association
records an observed relationship; a hypothesized cause belongs to a
candidate causal model. A household owns accounts, but the ownership edge is
not a causal arrow from a household variable to an account outcome. A synonym
link can support a search without authorizing the merger of two measurements.
An ontology reasoner checks implications of its formal premises; a causal
identification procedure has different premises and a different output.

Tan and colleagues' annotation survey separates generation, assessment and
utilization of annotations, and discusses hallucination, feedback from
synthetic data and manual or automatic validation.~\citep{exec17_ontology_annotation}
This separation suggests a staged bank workflow: an extractor proposes a
field and cites the passage; a typed validator checks identity, unit and
clock; a reviewer adjudicates meaning; a causal analysis uses the resulting
measurement under explicit assumptions. Agreement among model-generated
labels is evidence about agreement, not independent evidence about the
underlying construct. Generated respondents are especially unsuitable as
substitutes for observed customers in an effect estimate.

Schwitter, Bach and Klamm distinguish text as outcome, treatment and control
and illustrate substantive measurement decisions in sociological
applications. In one reported annotation example, zero-shot kappa is
\(0.63\pm0.01\), compared with \(0.78\pm0.01\) under three-shot prompting,
across three reruns. These are agreement results in the authors' setting,
not calibrated bank measurement errors or causal confidence intervals.
Their distinction between annotating real respondents and generating
respondent simulations is directly relevant here. Their illustrated noise
experiment discusses a tendency toward type-II error; differential
measurement need not behave that way.~\citep{exec17_ontology_sociology}

\section{High label accuracy can coexist with a false monetary conclusion}

Consider an exact authored population with randomized, equally sized offer
and control arms. The true adverse-event risk is \(p_0=p_1=1/5\). Write
\[
 s_a=P(L=1\mid Y=1,A=a),\qquad
 f_a=P(L=1\mid Y=0,A=a).
\]
The observed label risk is
\[
 q_a=P(L=1\mid A=a)=s_ap_a+f_a(1-p_a)
     =(s_a-f_a)p_a+f_a.
\]
This identity follows by partitioning on the true outcome. It does not
assume that an observed label is the true outcome. Suppose the control arm has
\((s_0,f_0)=(4/5,1/20)\), while the offered arm has
\((s_1,f_1)=(9/10,3/20)\). For instance, the intervention may change employees'
documentation and the terms customers use. The stipulated numbers imply
\[
 q_0=\frac15,\qquad q_1=\frac3{10},\qquad q_1-q_0=\frac1{10},
 \qquad p_1-p_0=0.
\]
Control accuracy is \(0.92\), offered-arm accuracy is \(0.86\), and pooled
accuracy is \(0.89\). Despite that apparently respectable aggregate accuracy,
the label comparison reports a ten-percentage-point adverse effect under a
true null. At a stipulated loss of USD~1,000 per true event, treating labels
as events invents USD~100 of additional loss per account. These are
illustrative exact values, not estimates from customers or model
calls.~\citep{exec17_ontology_oracle}

If valid arm-specific measurement rates were known and \(s_a\ne f_a\), the
algebraic correction
\[
 p_a=\frac{q_a-f_a}{s_a-f_a}
\]
would recover \(1/5\) in each arm in this example. The rates must describe the
actual measurement process and target population. A representative validation
sample would introduce estimation uncertainty; small \(s_a-f_a\) amplifies
that uncertainty; shifts in language or workflow can invalidate the
calibration. Prompt confidence is not a supplied value of \(s_a\) or \(f_a\).

Without measurement information the problem is one of identification, not
just classifier improvement. The same observed pair
\((q_0,q_1)=(1/5,3/10)\) also arises under perfect measurement and true risks
\((p_0,p_1)=(1/5,3/10)\). Thus the observed treatment-label law is compatible
with true effects zero and \(1/10\). More copies of that law cannot choose
between these worlds. A useful proposal can compare sensitivity ranges,
identify an additional validation measurement, or explicitly target the
effect on coded notes. It cannot recover a unique latent effect by renaming
the label.

This example explains why validation must be stratified by scientifically
relevant circumstances: treatment arm, language, document source and time
can alter error rates. The eventual sample design must justify how
adjudicated observations represent the target population. Such expensive
validation follows a promising prototype; the proposal can establish the
failure mechanism and candidate remedies using generated laws now.

\section{Entity, unit and time are part of the variable}

Aggregation can change the question while preserving every individual number.
Suppose two households have stipulated household effects of USD~\(10\) and
USD~\(-10\), but nine and one accounts respectively. The equal-household mean
is zero. Replicating each household's value once per account gives
\[
 \frac{9(10)+1(-10)}{9+1}=8.
\]
This is an account-weighted \emph{household} effect, not proof of an
account-level causal mechanism. Ten rows still refer to two households.
The representation must retain entity identifiers, grouping links,
the level of treatment assignment and the target weighting. A precision
calculation cannot treat repeated household records as independent households.
\citep{exec17_ontology_oracle}

Units require more than a monetary label. Monthly debt service of
USD~1,500 and monthly income of USD~3,000 imply a debt-service ratio of \(1/2\).
Dividing the same monthly obligation by annual income of USD~36,000 gives
\(1/24\). Converting annual income to monthly income restores \(1/2\).
The conversion is justified by an explicit common-period definition in this
example. It would not reconcile gross with net income, household with
individual income, a stock with a flow, or average with end-period balances.
Those are different constructs or aggregation rules. Currency conversion
would additionally require the rate and valuation date.

An observation's event time and usable time also differ. A fact about day
one that becomes available on day five cannot inform a day-three decision.
For a derived value computed from inputs with availability times two and
five, availability cannot precede five. A declared time of four fails even
when the input events occurred earlier. Actual production delay may make
availability later still. These examples define a historical information
constraint; they do not prohibit outcomes from becoming observable after
the treatment decision.~\citep{exec17_ontology_oracle}

The local proposal implementation already checks declared measured/latent
types, derived lineage, event-before-availability order, input availability
and information at a decision. Its restricted static-effect projection
retains variable identity, role and unit, keeps full declarations as evidence
and refuses a mismatch between outcome and loss units because it has no
utility transformation. These are useful declaration checks. They neither
audit individual rows nor supply a verified extraction, household linkage or
currency-conversion service.~\citep{exec17_ontology_sources}
The proposed next layer must carry record-specific clocks and provenance
through the conversion, and preserve an explicit unresolved state when
availability or identity cannot be established.

\section{Missing notes and imputed values need an observation model}

Even perfect annotation of the notes that exist can mislead. Let
\(R_i=1\) indicate that a note is observed. In another exact authored law, the
true adverse-event risk is \(1/5\) in both randomized arms. All control notes
are observed. In the offered arm, all adverse events have a note, but only
one quarter of non-events do. Then
\[
 P(R=1\mid A=1)=\frac15+\frac45\frac14=\frac25,
 \qquad
 P(Y=1\mid A=1,R=1)=\frac{1/5}{2/5}=\frac12.
\]
Comparing observed notes gives an apparent contrast \(1/2-1/5=3/10\)
under a true null. Selection depends jointly on treatment and outcome.
Annotation accuracy among observed notes cannot repair the missing
population.~\citep{exec17_ontology_oracle}

Known positive observation probabilities would permit the illustrated
weighted population calculation to recover \(1/5\); those probabilities are
stipulated in this reference, not estimated from notes with unobserved
outcomes. If the relevant selection law is unknown, collecting representative
outcomes or adding justified assumptions is necessary. An assistant can expose
the missingness mechanism, propose a sensitivity analysis and keep a value
marked missing. Filling it with a plausible sentence does not reveal what
the customer would have reported.

This caution is consistent with the primary quantitative-knowledge study
reviewed earlier. Its imputation comparisons report that KNN and random
forest methods outperform the examined LLM alternatives in the stated
settings; ambiguous descriptions and units are discussed as contributors to
error. Recognition of familiar datasets also raises contamination concerns.
\citep{v12_la_2402_07770} These findings support retaining ordinary imputation
and dictionary baselines. They do not imply that all LLM extraction fails,
or that improved prediction identifies the causal effect when observation
depends on an unobserved outcome.

The ontology should represent missing, not applicable, not collected,
redacted, delayed and model-imputed values distinctly. It should retain the
original observation indicator even if an imputation is created. A changed
imputation model creates a new derived measurement revision; it does not
rewrite the original record or convert a latent quantity into an observed
fact. Inference about imputed values needs a separate justification of the
imputation and sampling laws.

\section{Benchmark evidence and a concrete package decision}

OntoURL makes the distinction between understanding, reasoning and generation
testable through fifteen ontology tasks. Its construction uses rule
extraction, a reasoner for reasoning tasks, distractors and quality review.
The inspected preprint reports strong performance on some recognition tasks
and weak performance on several structured generation tasks. For example,
its Qwen2.5-72B row gives class-definition and class-relation understanding
(U1 and U2) accuracies of \(89.1\%\) and \(92.6\%\), and rule-based and
description-logic reasoning (R4 and R5) accuracies of \(75.6\%\) and \(68.4\%\).
Class-hierarchy, property-relation and constraint construction (L2--L4) have
Triple-F1 scores of \(0.1\), \(0.0\) and \(1.0\); ontology alignment (L5)
has Tuple-F1 of \(21.6\), on the table's percentage scale.
These are different tasks and metrics, not a causal-graph accuracy
scale.~\citep{exec17_ontology_ontourl}

The same version contains a material reporting inconsistency: Table~2 gives
L1 BERTScores around 79--83, while nearby prose says they are typically below
10 and Table~3 uses another apparent scale. No aggregate comparison should
silently reconcile these numbers. The three PhD-level human participants
answer thirty sampled questions per task, while the table's model results
cover the full dataset. This also limits a simple human-versus-model
ranking. Prompting effects vary across tasks, and the authors acknowledge
restricted domains, incomplete task coverage and an English-only scope.
The usable lesson is to evaluate the actual task and output semantics;
the disputed scale and sampling differences require separate resolution.

OntoLearner offers a more concrete reuse opportunity: a Python framework
with task-specific data and evaluation paths. Its paper reports
domain-dependent error patterns and calibration behavior, and itself argues
that F1 alone misses important failure modes.~\citep{v12_la_2607_01977}
The inspected revision's document loader retains document identifiers and
extraction methods, its evaluator dispatches among four task metrics, and
its non-taxonomic metric only treats three declared relation names as
symmetric. In particular, the code does not symmetrize every relation:
reversing a causal triple is different from reversing a
\texttt{sameas} triple.~\citep{exec17_ontology_sources}

There is a specific reason to adapt its evaluation interface before any
causal use. Its text-to-ontology metric concatenates the document identifier
and extracted value, then greedily matches token sets using a Jaccard
threshold of \(0.8\). Take two different document identifiers and the
same nine-word phrase, \emph{monthly household income before tax during the
reporting period}. The token intersection has size nine and the union
eleven, so similarity is \(9/11>0.8\) even though the document link is wrong.
This is an exact authored counterexample to the inspected scoring
definition, not an executed upstream-package result.
\citep{exec17_ontology_sources,exec17_ontology_oracle}

A causal observation contract should first require exact document identity,
then compare meaning within that document; source-span identity, units and
clocks need their own checks. Lexical similarity can remain a useful
diagnostic after those requirements. Similarly, a taxonomy metric can measure
hierarchical agreement without validating an intervention graph, and
deduplication can alter denominators if the intended unit is an observation
rather than a unique term.

The package decision is to use OntoLearner as a reference and candidate
benchmark component, with an adapter proposed for typed observations and
separate semantic scoring. Direct adoption of the inspected metrics as a
causal acceptance gate is rejected. The inspected downloads do not pin
dataset revisions; enabled ontology-import failures can be logged while
processing continues; missing dataset splits can also continue with warnings.
A reproducible evaluation must freeze required dependencies and splits and
explicitly fail or mark incomplete required missing inputs. The base
learner's formatting path does not establish all concrete subclass behavior.
Its MIT license permits relevant code reuse subject to the notice, but does
not establish rights to each external ontology, dataset, model or dependency.
Full runtime qualification remains a separate proposed step.
\citep{exec17_ontology_sources}

\section{A bounded public and generated comparison before bank validation}

The proposed prototype takes a document or record, a vocabulary revision, a
business question and its decision clock. It returns candidate observations
with the original record and span, entity and aggregation level, construct and
operational definition, value domain and unit, event and availability clocks,
transformation lineage, extraction revision, and an explicit evidence status.
Candidate causal roles and relations are separate objects linked to the
model in which they are proposed. Unresolved meaning remains unresolved.
An assistant's numeric confidence may be stored as a model output, but is
not interpreted as a calibrated measurement-error probability.
\citep{exec17_ontology_contract}

The initial comparison should contain a deterministic dictionary/schema
baseline, an ontology-assisted mapping alternative, and a constrained
language-model extractor if a later permitted model trial is warranted.
The six exact families above supply tractable semantic reference cases.
Additional generated variants should include ambiguous entities, negation,
incompatible definitions, conflicting passages, renamed identifiers and
unavailable facts. Positive cases establish what a valid conversion must
retain; negative cases establish the appropriate refusal or unresolved
output. A case with observationally equivalent causal worlds should reward
preserved ambiguity, not a lucky choice of the generator's hidden world.

Public ontology benchmarks can test recognition and formal relation tasks;
their distributions, rights and score definitions must be inspected before
use. Public treatment text and code identified by the sociology article are
possible comparison materials, whereas the article's onsite-only respondent
text is not a public dataset. None supplies bank causal truth. Generated
banking records can isolate known unit, clock, grouping and selection
mechanisms without pretending to be authentic bank observations. They also
allow newly constructed combinations that reduce straightforward lookup,
although they do not establish absence of contamination or real-world
generalization.

Before model comparison, freeze the case-generating rules, development and
evaluation split, document families, source revisions, permitted tools,
prompts, model identity, budgets and stopping rule. Keep paraphrases and
records from the same entity together when defining evaluation units. Freeze
the measurement mapping before effect evaluation. Report exact identity,
unit and clock preservation; typed relation direction; correctness of
acceptance, refusal and unresolved dispositions; performance on all requests
and on accepted outputs; and effects on the downstream target. Retain
malformed responses and missing documents in the denominator. Human
adjudication of meaning, if introduced later, needs an independent label
procedure and disagreements must remain visible.

Passing all stipulated semantic cases is a bounded engineering criterion,
not a guarantee of future causal validity. The comparison should separately
measure whether assistance saves review effort after retrieval, correction,
adjudication, retries and maintenance. A proposed economic calculation can
compare avoided review and error loss with these additional costs; it cannot
assign actual bank savings before the relevant error frequencies, prices
and consequences are measured. The next investment decision is whether a
permitted prototype demonstrates enough benefit to justify expensive
representative measurement validation.

The remaining research is specific. Historical availability and entity
matching need evidence for actual records; latent constructs need measurement
assumptions and representative validation; selection and imputation need a
defensible observation law; causal sufficiency needs additional scientific
assumptions or information. Better extraction can reduce some implementation
errors, but observationally equivalent measurement mechanisms require new
information or an explicitly broader answer. The proposed program can develop
and compare these remedies now without claiming that an ontology, benchmark
score or language model has settled them. The following studies ask how to
handle the evidence used to define those measurements and restrict their
causal models: how to distinguish textual knowledge from observations, retain
expert disagreement, and use advice without treating repeated agreement as
independent confirmation.
\label{exec17:ontology-end}
